\documentclass{article}
\usepackage[utf8]{inputenc}
\usepackage[T1]{fontenc}
\usepackage[slovene]{babel}

\title{Izdelava urnikov}
\author{Nejc Širovnik}
\date{14. 10. 2019}

\begin{document}

\maketitle

\begin{abstract}
V članku so predstavljeni nekateri postopki za avtomatsko izdelavo urnikov. Opisani so
potrebni vhodni podatki, naštete so omejitve, ki jih je potrebno upoštevati pri izdelavi
urnikov, predstavljena je tudi metoda merjenja kakovosti izdelanega urnika.
\end{abstract}

\section{Program za izdelavo urnika}
Spominjam se osnovnošolskih in srednješolskih dni, ko smo vsi z nestrpnostjo pričakovali
začetek novega šolskega leta...

\subsection{Vhodni podatki}
Če želimo izdelati urnik, v prvi vrsti potrebujemo podatke...

\subsection{Omejitve}
Programa za izdelavo urnika načeloma ni težko narediti...

\subsubsection{Stroge omejitve}
Verjetno še niste videli profesorja...

\subsubsection{Šibke omejitve}
V splošnem ne zahtevamo, da urnik zadošča vsem šibkim omejitvam...

\subsection{Merjenje kakovosti}
Kakovost urnika se meri z zadoščanjem šibkim omejitvam...

\section{Avtomatizirana izdelava urnika}
V 1. poglavju smo izvedeli, da ročna izdela urnikov ni tako preprosta...

\subsection{Požrešna metoda}
Požrešna metoda je ena izmed strategij...

\subsection{Optimizacijske metode}
Optimizacijske metode se že desetletja uporabljajo...

\subsubsection{Lokalno vzpenjanje}
Lokalno vzpenjanje temelji na opazkah...

\subsubsection{Iskanje s tabuji}
Že samo ime spominja na prepovedi...

\subsection{Končni algoritem}
Za konec si oglejmo, kako je videti optimizacijski del...

\subsection{Zaključek}
Sledili smo receptu in s tem izdelali okusen obrok...

\end{document}